\section{Cross-Dataset Hyperparameter Generalization}
\label{sec:tau_generalization}

The HCPC-v2 retriever exposes one tunable threshold $\tau$ that
gates whether a passage is refined. The published HCPC-v2 numbers
use $\tau$ tuned on a 50-query SQuAD held-out subset. This section
asks whether that choice of $\tau$ transfers to other datasets.

\paragraph{Protocol.}
For each dataset $\mathcal{D} \in \{\text{SQuAD}, \text{PubMedQA},
\text{HotpotQA}, \text{NaturalQuestions}, \text{TriviaQA}\}$ and
each candidate $\tau \in \{0.30, 0.40, 0.50, 0.60, 0.70\}$, we run
HCPC-v2 at $n{=}300$ queries and report
\(\mathrm{recovery}(\mathcal{D},\tau) =
(\overline{\phi}_{\text{ccs-gate}} - \overline{\phi}_{\text{HCPC-v1}})
\,/\,(\overline{\phi}_{\text{baseline}} - \overline{\phi}_{\text{HCPC-v1}})\).
Tuning ``on $\mathcal{D}$'' means picking the $\tau$ that
maximizes $\mathrm{recovery}(\mathcal{D},\tau)$. Evaluating ``on
$\mathcal{D}'$'' means using that $\tau$ to score
$\mathrm{recovery}(\mathcal{D}',\tau)$ on the held-out dataset.
Total: $5\times 5 \times 300 = 7{,}500$ evaluations.

\paragraph{Result.}
For each tuning dataset, we report the diagonal recovery
($\mathcal{D}=\mathcal{D}'$) versus the mean off-diagonal recovery
($\mathcal{D}\ne\mathcal{D}'$). A pre-registered flag fires when
the gap exceeds $0.03$.

\begin{table}[H]
\centering
\small
\caption{Cross-dataset $\tau$ generalization. ``Flag'' is set when
\emph{(diagonal recovery $-$ off-diagonal mean recovery) $> 0.03$};
flagged datasets must be acknowledged honestly in any HCPC-v2
result that uses their tuned $\tau$.}
\label{tab:tau-generalization}
\begin{tabular}{lcccc}
\toprule
Tune dataset & Diagonal recovery & Off-diag mean recovery &
Diag $-$ Off-diag & Flag \\
\midrule
PubMedQA           & $1.452$ & $0.451$ & $1.002$ & \textbf{Yes} \\
NaturalQuestions   & $1.401$ & $0.211$ & $1.191$ & \textbf{Yes} \\
SQuAD              & $0.823$ & $-0.141$ & $0.963$ & \textbf{Yes} \\
TriviaQA           & $0.464$ & $0.698$ & $-0.234$ & No \\
HotpotQA           & $-0.112$ & $0.093$ & $-0.205$ & No \\
\bottomrule
\end{tabular}
\end{table}

\paragraph{Interpretation.}
Three datasets --- the SQuAD on which the published $\tau$ was
tuned, plus PubMedQA and NaturalQuestions --- exhibit substantial
in-distribution recovery that does not transfer to other datasets.
Two datasets (TriviaQA, HotpotQA) generalize better than they
in-distribution-tune.

For SQuAD specifically: tuning $\tau$ on a 50-query SQuAD held-out
yields a recovery of $0.823$ on the rest of SQuAD --- but when that
same $\tau$ is evaluated on the other four datasets the mean
recovery is \emph{negative} ($-0.141$). The headline HCPC-v2
recovery number on SQuAD is therefore amplified by $\tau$ tuning,
not just by the method.

\paragraph{What this means generally.}
Threshold-based RAG interventions are vulnerable to in-distribution
$\tau$ tuning even when researchers report tuning on a
``small held-out subset.'' The convention of reporting a single
$\tau$ per system without disclosing which subset it was tuned on,
or how the choice transfers across datasets, is an evaluation
practice that systematically inflates within-dataset numbers
relative to honest off-distribution numbers.
For the rest of this paper we report all HCPC-v2 results at the
\emph{cross-dataset off-diagonal} $\tau$ rather than the
in-distribution diagonal $\tau$, which moves the SQuAD recovery
estimate from $0.823$ to $-0.141$. The honest summary: the
``HCPC-v2 recovers faithfulness'' claim depends on which $\tau$ is
chosen, and the chosen $\tau$ in the published work was tuned on
the very dataset on which the claim was reported.
